% !TEX encoding = UTF-8
% !TEX spellcheck = fr_FR
% !TEX program = lualatex
%
% titrages-technique.tex — les calculs, les choix algorithmiques et les vérifications du
% paquet titrages (Python, Lua, LaTeX).
%
% Compilation, depuis ce dossier, le paquet LaTeX titrages étant pris dans ../../latex :
%   uv run python verifications.py            (refait tables/, une minute)
%   TEXINPUTS=../../latex: lualatex titrages-technique
%   TEXINPUTS=../../latex: lualatex titrages-technique
%
% Sources :
%   H. Kalka, Polyprotic Acids and Beyond — An Algebraic Approach, Chemistry 2021, 3, 454-508.
%   T. Brière, Une méthode simple de résolution des problèmes de pH et de titrages
%     acido-basiques sans approximations, BUP 841, 2002.
%   Y. Cortial, Courbe de dosage acidobasique et méthode des tangentes, BUP 734, 1991.
%   F. M. M. Morel, J. G. Hering, Principles and Applications of Aquatic Chemistry, Wiley, 1993
%     (la méthode du tableau, les composantes).
%   D. Harvey, Analytical Chemistry 2.1, LibreTexts, § 16.10 à 16.12 (constantes).
%
% Copyright (c) 2026 antnardo — licence MIT (voir LICENSE).

\documentclass[a4paper,11pt]{article}
\usepackage{fontspec}
\usepackage[french]{babel}
\usepackage[margin=2.2cm]{geometry}
\usepackage{amsmath,amssymb}
\usepackage{mathtools}
\usepackage[version=4]{mhchem}
\usepackage{siunitx}
\sisetup{locale=FR, exponent-product=\cdot}
\usepackage{booktabs,tabularx}
\usepackage{enumitem}
\usepackage[hypcap=false, justification=raggedright, singlelinecheck=false]{caption}
\usepackage{xcolor}
\usepackage[most]{tcolorbox}
\usepackage{fontawesome5}
\usepackage{pifont}
\usepackage{xspace}
\usepackage[hidelinks]{hyperref}
\usepackage{titrages}
\emergencystretch=3em
\pgfplotsset{compat=1.18}

% --- les boîtes, à l'image des fiches dont ce document reprend le style -----------------
\tcbset{
  rmq/.style={enhanced, breakable, boxrule=0pt, frame hidden, arc=2pt, outer arc=2pt,
    left=15pt, right=3pt, top=4pt, bottom=4pt, boxsep=0pt, before skip=4pt, after skip=4pt,
    fontupper=\small},
  bleu/.style={interior style={left color=blue!10, right color=black!5}},
  rouge/.style={interior style={left color=red!20, right color=red!5}},
  vert/.style={interior style={left color=green!40!black!30, right color=green!40!black!10}},
}
\newcommand\logo[1]{\hspace{-12pt}\makebox[12pt][l]{#1}}
\newcommand\rmqphys[1]{\begin{tcolorbox}[rmq, bleu]\logo{\faLightbulb}#1\end{tcolorbox}}
\newcommand\rmqplus[1]{\begin{tcolorbox}[rmq, bleu]\logo{\ding{40}}#1\end{tcolorbox}}
\newcommand\rmqcons[1]{\begin{tcolorbox}[rmq, rouge, fontupper=\small\itshape]\logo{\ding{43}}#1\end{tcolorbox}}
\newcommand\rmqattention[1]{\begin{tcolorbox}[rmq, vert]\logo{\faExclamationTriangle}#1\end{tcolorbox}}
\newcommand\rmxphys[1]{\begin{tcolorbox}[rmq, bleu, left=3pt]\begin{itemize}[label=\faLightbulb, leftmargin=*, topsep=0pt, itemsep=2pt]#1\end{itemize}\end{tcolorbox}}
\newcommand\rmxcons[1]{\begin{tcolorbox}[rmq, rouge, left=3pt, fontupper=\small\itshape]\begin{itemize}[label=\ding{43}, leftmargin=*, topsep=0pt, itemsep=2pt]#1\end{itemize}\end{tcolorbox}}
\newcommand\asavoir{\par\noindent\textbf{\textsf{À savoir.}}\xspace}
\newcommand\gras[1]{\textbf{#1}}
\newcommand\rsltm[1]{\boxed{#1}}
\newcommand\rslt[1]{\fbox{#1}}
\newcommand\captionout[1]{\captionof{figure}{#1}}

% --- notations ---------------------------------------------------------------------------
\newcommand\pH{\ensuremath{\mathrm{pH}}\xspace}
\newcommand\pL{\ensuremath{\mathrm{pL}}\xspace}
\newcommand\pM{\ensuremath{\mathrm{pM}}\xspace}
\newcommand\pP{\ensuremath{\mathrm{pP}}\xspace}
\newcommand\pKa{\ensuremath{\mathrm{p}K_\mathrm{a}}\xspace}
\newcommand\pKs{\ensuremath{\mathrm{p}K_\mathrm{s}}\xspace}
\newcommand\pKe{\ensuremath{\mathrm{p}K_\mathrm{e}}\xspace}
\newcommand\Ke{\ensuremath{K_\mathrm{e}}\xspace}
\newcommand\co[1]{[\ce{#1}]}
\newcommand\ddp[2]{\frac{\mathrm{d}#1}{\mathrm{d}#2}}
\newcommand\pdp[2]{\frac{\partial#1}{\partial#2}}
\newcommand\pare[1]{\left(#1\right)}
\newcommand\et{\quad\text{et}\quad}
\newcommand\Lra{\Longleftrightarrow}
\newcommand\code[1]{\texttt{#1}}
\newcommand\Rb{R_\text{bécher}}
\newcommand\Ru{R_\text{burette}}

\title{Le paquet \code{titrages}\\[4pt] \large Calculs, choix algorithmiques et vérifications}
\author{antnardo}
\date{19 octobre 2026}

\begin{document}
\maketitle

Ce document accompagne le paquet \code{titrages} — un paquet Python et un paquet \LaTeX{} qui calculent la courbe de titrage d'un bécher quelconque par une burette quelconque, sans aucune approximation.
Il donne tous les calculs sur lesquels reposent les deux moteurs, explique chacun des choix algorithmiques, et rassemble les preuves de leur fonctionnement : les démonstrations d'unicité et de monotonie qui fondent les algorithmes, et les comparaisons à des calculs indépendants, en double précision ou en précision arbitraire, dont les tables donnent les écarts mesurés.

\rmxcons{
\item Les solutions sont idéales : l'activité est la concentration.
La température n'intervient que par \pKe et par la pente de Nernst.
\item Les concentrations sont en \unit{mol/L}, les volumes dans une unité quelconque, la même partout.
\item Le code est dans \code{src/titrages/} (Python) et \code{latex/titrages.lua} (Lua) ; chaque section nomme le module qu'elle décrit.
Les tables de la dernière partie sont écrites par \code{docs/technique/verifications.py}, qui les recalcule.
}

\tableofcontents

% =========================================================================================
\section{Les fractions d'un polyacide}

\subsection{Notations}
\begin{itemize}
\item $h=\co{H3O+}$, $w=\co{HO-}$, et le produit ionique de l'eau $\Ke=hw$.
\item Un polyacide sous $N+1$ formes, de la plus acide (forme $0$) à la plus basique (forme $N$), de concentration totale $C$ ; la forme $j$ a cédé $j$ protons.
\item Ses fractions $x_j=a_j/C$, avec $a_j$ la concentration de la forme $j$.
\item Les constantes d'acidité successives $K_j$ et les constantes de dissociation cumulées $k_j=\prod_{i\le j}K_i$, constantes de \ce{H_{$N$}A = $j$ H+ + H_{$N-j$}A^{$-j$}} ; $k_0=1$.
\item La forme $q$ est celle sous laquelle on introduit l'espèce : \ce{NaHCO3} introduit \ce{HCO3-}, $q=1$.
\end{itemize}

\subsection{Les fractions, explicites}
La loi d'action des masses donne $k_j=h^ja_j/a_0$, et la conservation de la matière $\sum_jx_j=1$, d'où
\begin{align}
x_j(h)=\frac{k_j/h^j}{\sum_{m=0}^Nk_m/h^m}.
\label{eq:fractions}
\end{align}
Avec $\log(k_j/h^j)=\log k_j+j\,\pH$, on note $t_j=\log k_j+j\,\pH$ et $t_\text{max}=\max_mt_m$ ; alors
\begin{align}
\rsltm{x_j=\frac{10^{t_j-t_\text{max}}}{\sum_m10^{t_m-t_\text{max}}}}
\label{eq:lse}
\end{align}
\rmqphys{Le plus grand terme du dénominateur vaut exactement 1, les autres sont entre 0 et 1 : la somme est entre 1 et $N+1$, et aucune fraction n'est le quotient de deux nombres démesurés.
Écrite comme \eqref{eq:fractions}, la formule déborde : pour l'EDTA à $\pH=14$, $\log(k_6/h^6)=-\num{22,56}+84=\num{61,4}$, et $x_0=10^{-61}$.
C'est le procédé dit \og log-sum-exp\fg{} (\code{equilibre.fractions}, \code{fractions\_espece} en Lua).}

\subsection{Les moments}
Le nombre moyen de protons cédés, et le moment d'ordre 2,
\begin{align}
Y_1=\sum_{j=0}^Nj\,x_j \et Y_2=\sum_{j=0}^Nj^2x_j .
\end{align}
De $\ln x_j=\ln10\,t_j-\ln\sum_m10^{t_m}$, et de $\ddp{t_j}{\pH}=j$,
\begin{align}
\ddp{x_j}{\pH}=\ln10\,x_j\pare{j-Y_1}
\quad\Longrightarrow\quad
\ddp{Y_1}{\pH}=\ln10\sum_jj\,x_j\pare{j-Y_1}=\ln10\pare{Y_2-Y_1^2}.
\label{eq:dY1}
\end{align}
$Y_2-Y_1^2$ est la variance du nombre de protons cédés : $Y_1$ est \gras{strictement croissante} en \pH dès que deux formes coexistent.

\subsection{La variance, sans soustraction}
Dans un saut, une forme domine à $1-10^{-10}$ près : $Y_2$ et $Y_1^2$ coïncident sur dix chiffres, et leur différence n'en garde que six.
On l'écrit donc autrement :
\begin{align}
Y_2-Y_1^2
=\sum_jj^2x_j-\sum_{j,k}j\,k\,x_jx_k
=\frac12\sum_{j,k}x_jx_k\pare{j^2+k^2-2jk}
=\rsltm{\sum_{j<k}x_jx_k\pare{k-j}^2}
\label{eq:variance}
\end{align}
en utilisant $\sum_kx_k=1$ pour symétriser.
Tous les termes sont positifs : aucun chiffre ne se perd.
De même, l'écart de $Y_1$ à la forme introduite s'écrit $Y_1-q=\sum_jx_j\pare{j-q}$, sans soustraire deux nombres voisins (\code{equilibre.variance}, \code{equilibre.ecart\_moyen}).

% =========================================================================================
\section{Le pH d'une solution}

\subsection{L'électroneutralité}
Une solution contient des espèces faibles $i$, de concentrations $C_i$, introduites sous les formes $q_i$ ; un acide fort $C_a$ et une base forte $C_b$ ; leurs contre-ions.
Une espèce introduite sous la forme $q_i$ apporte un contre-ion qui compense sa charge : ce qui compte est l'écart des protons cédés à ceux qu'on a ôtés en l'introduisant, $C_i(Y_{1,i}-q_i)$.
La conservation de la charge s'écrit alors
\begin{align}
\rsltm{R(\pH)=\sum_iC_i\pare{Y_{1,i}-q_i}-\pare{C_b-C_a}+\delta}
\qquad\text{avec}\qquad\delta=w-h,
\label{eq:R}
\end{align}
et le \pH de la solution est celui qui annule $R$ (\code{Solution.residu}).

\subsection{Unicité}
D'après \eqref{eq:dY1} et $\ddp\delta\pH=\ln10\pare{w+h}$,
\begin{align}
\ddp R\pH=\ln10\pare{\sum_iC_i\pare{Y_{2,i}-Y_{1,i}^2}+h+w}>0 :
\label{eq:dR}
\end{align}
$R$ est strictement croissante, et s'annule au plus une fois.
Elle s'annule au moins une fois sur $[-6;\pKe+6]$ : en $\pH=-6$, $h=\qty{e6}{mol/L}$ l'emporte sur toutes les concentrations, et $R<0$ ; en $\pH=\pKe+6$, $w=\qty{e6}{mol/L}$, et $R>0$.
\rmqphys{$\ddp R\pH$ est le pouvoir tampon de la solution, $\ddp{C_b}{\pH}$ (\code{Solution.pouvoir\_tampon}) : les deux notions sont la même dérivée.}

\subsection{La dichotomie}
Le \pH se trouve par dichotomie sur $[-6;\pKe+6]$ : on coupe l'intervalle en deux, on garde la moitié où $R$ change de signe, 64 fois.
L'intervalle final mesure $(\pKe+12)/2^{64}\approx\num{1.4e-18}$, sous la précision de la machine au voisinage de 7 ($\num{8.9e-16}$) : le résultat est le \pH exact en double précision.
\rmxcons{
\item Pourquoi pas Newton : $\ddp R\pH$ varie de plusieurs ordres de grandeur entre les zones tampon et les sauts.
Un pas de Newton parti d'une zone où la pente est presque nulle part à l'infini ; la dichotomie ne regarde que le signe, et ne demande ni réglage ni point de départ.
\item Les 64 coupes sont les mêmes pour tous les points : elles se font sur un tableau de volumes d'un coup (\code{equilibre.racine\_croissante}).
}

% =========================================================================================
\section{Le titrage : le sens direct}

\subsection{Le mélange}
Verser $V$ de burette dans $V_0$ de bécher, c'est mélanger les deux solutions : chaque concentration du mélange est la moyenne des deux pondérée par les volumes, $C^\text{mél}=\pare{V_0C^\text{b}+VC^\text{u}}/\pare{V_0+V}$.
À \pH fixé, les $Y_{1,i}$ ne dépendent pas des concentrations : le premier terme de \eqref{eq:R} est linéaire en concentrations, et $\delta$ ne dépend que du \pH.
L'électroneutralité du mélange, multipliée par $V_0+V$, se sépare donc en une part du bécher et une part de la burette, chacune calculée avec ses concentrations d'avant le mélange :
\begin{align}
(V_0+V)\,R^\text{mél}(\pH)=V_0\,\Rb(\pH)+V\,\Ru(\pH)=0
\quad\Longrightarrow\quad
\rsltm{V(\pH)=-V_0\,\frac{\Rb(\pH)}{\Ru(\pH)}}
\label{eq:V}
\end{align}
\rmxphys{
\item La dilution ne demande aucun traitement : il suffit de ne pas diluer les concentrations, et de laisser $V_0$ et $V$ pondérer les deux résidus.
\item La courbe a un sens là où $V\ge0$ : entre le \pH du bécher ($\Rb=0$, $V=0$) et celui de la burette ($\Ru=0$, $V\to\infty$), là où les deux résidus sont de signes opposés.
\item Ni polynôme à résoudre, ni réaction prépondérante : tout est explicite à \pH fixé (\code{Titrage.volume}).
}

\rmqattention{L'article de Kalka (2021) trace $n_0(h)=Y_1+\delta/C$ comme une courbe de titrage, ce qui revient à négliger la dilution.
Pour l'acide éthanoïque et la soude à \qty{0.1}{mol/L}, l'écart atteint $\log3\approx\num{0,48}$ unité de \pH à deux fois l'équivalence (\code{tests/test\_dilution.py}).}

\gras{Deux cas particuliers.}
\begin{itemize}
\item Le \pH imposé à une solution d'une seule espèce donne sa concentration, sans résolution : $R=0$ avec $C_a=C_b=0$ s'écrit $\rsltm{C(h)=-\delta(h)/\pare{Y_1(h)-q}}$ (\code{concentration\_pour\_pH}), le sens direct de la courbe $\pH=f(\mathrm pC)$.
\item Un bécher d'acide faible ($\Rb=C_TY_1+\delta$) et une burette de base forte ($\Ru=-C_B+\delta$) donnent, avec l'abscisse normalisée $n=C_BV/(C_TV_0)$,
\begin{align}
n(h)=\frac{Y_1(h)+\delta(h)/C_T}{1-\delta(h)/C_B} :
\end{align}
le numérateur seul est la courbe sans dilution, et le dénominateur, la correction de dilution.
\end{itemize}

\gras{Sans dilution.} L'option \code{dilution=False} décrit l'ajout d'un réactif sans eau : le terme $\delta$ du mélange n'est plus porté que par le bécher, et $\Ru$ se réduit à sa partie $\sum_iC_i(Y_{1,i}-q_i)-(C_b-C_a)$ (\code{Solution.charge}).

\subsection{La pente}
En dérivant $V_0\Rb+V\Ru=0$ par rapport au \pH,
\begin{align}
V_0\Rb'+V'\Ru+V\Ru'=0
\quad\Longrightarrow\quad
\rsltm{\ddp V\pH=-\frac{V_0\Rb'+V\Ru'}{\Ru}}
\label{eq:dV}
\end{align}
Le numérateur est une somme de deux termes positifs \eqref{eq:dR}, à poids positifs : rien ne s'y soustrait.
Le quotient $V_0\pare{\Rb\Ru'-\Rb'\Ru}/\Ru^2$, équivalent sur le papier, est au contraire une différence de deux termes presque égaux au cœur d'un saut.
La pente de la courbe est l'inverse, $\ddp\pH V$.

\subsection{L'échantillonnage}
La courbe est une courbe paramétrée par le \pH.
On calcule $V(\pH)$ sur une grille de $40n$ \pH réguliers, entre le \pH du bécher et celui du volume maximal, puis on choisit $n$ points à \gras{longueur d'arc constante} dans le plan normalisé $\pare{V/V_\text{max},\pH/\Delta\pH}$ : les points se resserrent dans les sauts, où la courbe tourne, et s'espacent sur les plateaux.
Chaque point retenu est recalculé exactement, et non interpolé (\code{Titrage.courbe}).

\begin{center}
\titragedefinir{ethanoique}{volume = 10, becher = { espece = { pKa = 4.76, C = 0.1 } }, burette = { base forte = 0.1 }, volume max = 20}
\begin{tikzpicture}
\begin{axis}[width=.62\textwidth, height=5.2cm, xlabel={$V$ (mL)}, ylabel=\pH, grid=both, ymin=0, ymax=14]
  \addplottitrage[thick, mark=*, mark size=.6pt, only marks]{ethanoique}
  \addplottitragesauts[red, mark=*]{ethanoique}
\end{axis}
\end{tikzpicture}
\captionout{\qty{10}{mL} d'acide éthanoïque à \qty{0.1}{mol/L} par la soude à \qty{0.1}{mol/L}, 300 points à longueur d'arc constante : ils se resserrent dans le saut. Point d'inflexion à \qty{\titragesaut{ethanoique}{1}}{mL}, \pH \num{\titragesaut[pH]{ethanoique}{1}}.}
\end{center}

% =========================================================================================
\section{Le titrage : le sens inverse}

Le \pH après un volume $V$ donné annule
\begin{align}
g(\pH)=V_0\,\Rb(\pH)+V\,\Ru(\pH),
\end{align}
somme de deux fonctions strictement croissantes \eqref{eq:dR} à poids positifs : $g$ est strictement croissante, et la dichotomie de la section précédente s'applique telle quelle, sur le même intervalle (\code{Titrage.pH}).
\rmqcons{On pourrait aussi interpoler la courbe directe ; mais c'est exact à la finesse de la grille près, quand la dichotomie l'est à la précision de la machine, pour le même coût.}

% =========================================================================================
\section{Les points d'inflexion}

\subsection{Ce qu'on cherche}
Un point d'inflexion de la courbe $\pH(V)$ est un extremum de sa pente $\ddp\pH V$.
On le cherche sur la courbe directe, comme minimum de $\left|\ddp V\pH\right|$ : le changement de paramètre ne déplace pas l'extremum, puisque
\begin{align}
\ddp{}{\pH}\pare{\ddp V\pH}=-\pare{\ddp V\pH}^3\frac{\mathrm d^2\pH}{\mathrm dV^2}
\end{align}
s'annule en même temps que la dérivée seconde de $\pH(V)$, la pente $\ddp V\pH$ étant finie et non nulle.

\subsection{Pourquoi en \pH, et non en volume}
Un saut tient parfois dans un intervalle de volume minuscule : celui de \ce{Fe^2+} par \ce{MnO4-} occupe moins de \qty{e-9}{mL}.
Aucune grille en $V$ ne le résout ; en \pH (ou en potentiel), il s'étale sur des unités, et $\left|\ddp V\pH\right|$ y est une fonction lisse.

\subsection{L'affinage : le zéro d'une dérivée, et non la comparaison de deux valeurs}
On repère les minima locaux de $f=\left|\ddp V\pH\right|$ sur une grille de 4000 \pH, puis on affine chacun entre ses deux voisins.
\begin{itemize}
\item Une section dorée compare $f(x_1)$ et $f(x_2)$.
Près d'un minimum, $f(x)=f(x^\star)+\tfrac12f''(x^\star)(x-x^\star)^2$ : deux valeurs ne s'y distinguent qu'au second ordre, et la précision sur $x^\star$ plafonne à $\sqrt\varepsilon\approx10^{-8}$ en relatif.
\item On cherche plutôt le zéro de $(\ln f)'$, estimée par différence centrée, $\ln f(x+\eta)-\ln f(x-\eta)$, de pas $\eta=10^{-6}$.
Une dérivée s'annule au premier ordre, et la dichotomie sur son signe atteint la précision de la machine (\code{equilibre.minimum\_lisse}).
\end{itemize}
\rmqphys{Les trois précautions — chercher dans le paramètre et non en volume, annuler une dérivée plutôt que comparer des valeurs, écrire la variance et la pente sans soustraction — se mesurent sur \ce{Fe^2+} par \ce{MnO4-}, face à mpmath à 60 chiffres (\qty{1,37288297818}{V}) : sans la première, \qty{1,3716}{V} ; sans la deuxième, \qty{1,37289}{V} ; sans la troisième, \qty{1,37300}{V}.
Avec les trois, l'écart est de quelques $10^{-12}$\,V (table~\ref{tab:inflexions}).}

\rmqattention{Un point d'inflexion n'est pas l'équivalence.
Pour \ce{Fe^2+} par \ce{MnO4-}, \qty{1,3729}{V} contre $E(V_\text{éq})=\qty{1,3867}{V}$, alors que les volumes ne diffèrent que de \qty{6e-10}{mL}.
L'acide phosphorique a trois équivalences et deux points d'inflexion seulement.}

% =========================================================================================
\section{Les équivalences stœchiométriques}

Elles ne demandent aucune résolution : ce sont des décomptes de protons.
Le bécher offre des \og sites\fg : chaque acidité $j$ encore portée par une espèce introduite sous la forme $q$ ($j>q$) est un site acide de \pKa $\mathrm{p}K_{\mathrm{a},j}$, de taille $C\,V_0$ ; chaque proton déjà ôté ($j\le q$) est un site basique.
\begin{enumerate}
\item Un excès de base forte dans le bécher consomme d'abord les sites acides, par \pKa croissant, et en fait des sites basiques de même \pKa ; un excès d'acide fort fait l'inverse.
Ce qui en reste est un site de \pKa infini.
\item Une base forte titrante prend les sites acides par \pKa croissant ; un acide fort, les sites basiques par \pKa décroissant.
\item La $k$-ième équivalence est le volume qui apporte autant de réactif que les $k$ premiers sites en demandent ; deux sites de même \pKa ne font qu'une équivalence.
\end{enumerate}
Un titrage en retour en découle : l'acide éthanoïque à \qty{0.05}{mol/L} et la soude à \qty{0.1}{mol/L}, titrés par l'acide chlorhydrique, donnent deux équivalences, à 5 et \qty{10}{mL} pour \qty{10}{mL} de bécher (\code{Titrage.volumes\_equivalents}).
Le titrant faible n'a pas de définition aussi simple : la fonction le refuse.

% =========================================================================================
\section{La méthode des tangentes}

Elle s'exécute exactement sur la courbe exacte, pour en mesurer l'erreur plutôt que la supposer.
\begin{enumerate}
\item Au volume $V_1$ choisi avant le saut, la pente $p_1=\ddp\pH V$, par \eqref{eq:dV}.
\item Après le premier saut suivant $V_1$, le point $\pH_2$ de même pente : le zéro de $\left|\ddp\pH V\right|-|p_1|$, encadré sur une grille de 2000 \pH entre le saut et le \pH du volume maximal, puis affiné par dichotomie sur son signe.
\item La parallèle équidistante passe par le milieu $\pare{(V_1+V_2)/2,(\pH_1+\pH_2)/2}$ ; son intersection avec la courbe est le zéro de $\pH-\pH_m-p_1\pare{V(\pH)-V_m}$ entre $\pH_1$ et $\pH_2$.
\end{enumerate}
Pour l'acide éthanoïque par la soude, depuis \qty{7}{mL} : \qty{\titragetangentes{ethanoique}{7}}{mL}, \pH \num{\titragetangentes[pH]{ethanoique}{7}}.

% =========================================================================================
\section{La conductimétrie}

La loi de Kohlrausch, $\sigma=\sum_i\lambda_i^\circ c_i$, est linéaire en concentrations ; à \pH fixé, toutes le sont, et
\begin{align}
\rsltm{\sigma=\frac{V_0\,\sigma_\text{bécher}(\pH)+V\,\sigma_\text{burette}(\pH)}{V_0+V}+\lambda^\circ_{\ce{H3O+}}h+\lambda^\circ_{\ce{HO-}}w}
\end{align}
où $\sigma_\text{bécher}$ et $\sigma_\text{burette}$ sont les parts des ions de chaque solution, formes chargées et ions spectateurs (le \ce{Na+} de la soude), calculées avec leurs concentrations d'avant le mélange et au \pH du mélange.
L'unité du paquet est le \unit{S/m}, celle que donnent $\lambda^\circ$ en \unit{mS.m^2.mol^{-1}} et $c$ en \unit{mol/L} ; la correction de dilution multiplie par $(V_0+V)/V_0$.

% =========================================================================================
\section{L'oxydoréduction}

\subsection{La transposition}
Les électrons remplacent les protons, et $E/s$ le \pH, avec $s=RT\ln10/F$.
Un élément a ses formes de la plus réduite (0) à la plus oxydée ($N$) ; l'étape $j$ échange $n_j$ électrons et consomme $m_j$ protons à la réduction ; $e_j=\sum_{i\le j}n_i$ est le nombre d'électrons cédés de la forme 0 à la forme $j$.
Le milieu est tamponné : $E^{\circ\prime}_j=E^\circ_j-s\,\frac{m_j}{n_j}\pH$.
La loi de Nernst donne $a_j/a_{j-1}=10^{n_j(E-E^{\circ\prime}_j)/s}$, d'où, comme \eqref{eq:lse},
\begin{align}
x_j=\frac{10^{t_j}}{\sum_m10^{t_m}}
\qquad\text{avec}\qquad
t_j=-\sum_{i\le j}\frac{n_iE^{\circ\prime}_i}s+e_j\frac Es,
\end{align}
et $Y=\sum_je_jx_j$, de dérivée $\ddp YE=\frac{\ln10}s\sum_{j<k}x_jx_k(e_k-e_j)^2>0$, par le même calcul que \eqref{eq:dY1} et \eqref{eq:variance}.
Le bilan d'électrons remplace l'électroneutralité : $R(E)=\sum_iT_i\pare{Y_i(E)-e_{q_i}}$, strictement croissant, et
\begin{align}
\rsltm{V(E)=-V_0\,\frac{\Rb(E)}{\Ru(E)}}
\end{align}
comme en acide-base (\code{redox.py}).

\subsection{Le sens inverse}
Le potentiel à volume donné annule $g(E)=V_0\Rb(E)+V\Ru(E)$, croissante : dichotomie sur $[\qty{-3}{V};\qty{4}{V}]$/ENH, qui contient tous les couples en solution aqueuse, 64 coupes ; l'intervalle final mesure $\qty{7}{V}/2^{64}\approx\qty{4e-19}{V}$.
Un bécher qui ne contient qu'une forme par couple (du \ce{Fe^2+} pur) n'a pas de potentiel : $\Rb>0$ partout, et la dichotomie converge vers une borne.
On le détecte ainsi — un résultat à moins de \qty{1}{\micro V} d'une borne —, et la courbe commence alors au millième du volume tracé.

\subsection{Les couples dimériques}
Pour \ce{I2}/\ce{I-}, \ce{Cr2O7^2-}/\ce{Cr^3+} ou \ce{S4O6^2-}/\ce{S2O3^2-}, la loi de Nernst porte un carré, $o=\kappa r^2$ avec $\kappa=10^{n(E-E^{\circ\prime})/s}$, et la conservation de l'élément, comptée en unités, $T=r+2\kappa r^2$.
Sa racine positive s'écrit sans soustraction quand $\kappa T\ll1$ :
\begin{align}
r=\frac{2T}{1+\sqrt{1+8\kappa T}}.
\end{align}
La fraction oxydée dépend alors de $T$ : il n'existe plus de $Y(E)$, une simple dilution change le potentiel, et la séparation \eqref{eq:V} ne tient plus.
\begin{itemize}
\item \gras{Le sens inverse tient toujours} : $Y(E,T)$ reste croissante en $E$ à $T$ fixé, et chaque dimère s'évalue à sa concentration du mélange dans $g(E,V)$.
\item \gras{Un seul dimère} : notons $u$ la concentration de sa forme monomère (\ce{Cr^3+}).
À $E$ fixé, les couples monomères contribuent au bilan, multiplié par $V_0+V$, pour une quantité affine en $V$, $A(E)+B(E)\,V$.
Avec pour inconnue $W=(V_0+V)u$, le bilan d'électrons d'un dimère oxydé (l'autre cas s'en déduit en échangeant les deux formes) s'écrit $A+BV+e_1\pare{V_0T_\text{b}+VT_\text{u}-W}-V_0I_\text{b}-VI_\text{u}=0$ ($I$ : les électrons introduits), linéaire en $W$ et en $V$ :
\begin{align}
W=w_0+w_1V,\qquad w_0=\frac{A+V_0\pare{e_1T_\text{b}-I_\text{b}}}{e_1},\qquad w_1=\frac{B+e_1T_\text{u}-I_\text{u}}{e_1}.
\end{align}
La conservation de l'élément multipliée par $(V_0+V)^2$, $\pare{V_0T_\text{b}+VT_\text{u}-W}(V_0+V)=2\kappa W^2$, est alors du second degré en $V$ : avec $d_0=V_0T_\text{b}-w_0$ et $d_1=T_\text{u}-w_1$, $c_2V^2+c_1V+c_0=0$, où
\begin{align}
c_2=d_1-2\kappa w_1^2,\qquad c_1=d_0+V_0d_1-4\kappa w_0w_1,\qquad c_0=V_0d_0-2\kappa w_0^2.
\end{align}
$V(E)$ est la racine qui a un sens physique — $V\ge0$, $W\ge0$, $V_0T_\text{b}+VT_\text{u}-W\ge0$, la plus petite si deux conviennent —, calculée par $\mathcal Q=-\frac12\pare{c_1+\operatorname{sgn}(c_1)\sqrt{c_1^2-4c_2c_0}}$, puis $V=\mathcal Q/c_2$ et $V=c_0/\mathcal Q$ : la formule scolaire perd tous ses chiffres quand $c_1^2\gg|4c_2c_0|$, ce qui arrive sur les plateaux, où $\kappa$ vaut $10^{\pm30}$.
\item \gras{Deux dimères ou plus} : $V(E)$ par dichotomie sur $\log V$ entre $10^{-12}V_0$ et $10^6V_0$, sur le signe de $g(E,V)$.
\end{itemize}

\gras{La pente d'un dimère.} Elle se prend par dérivation implicite, $\ddp EV=-\pdp gV\big/\pdp gE$, au point où $g$ est exactement nulle.
Avec $Y=e_1\pare{1-r/T}$ pour un dimère oxydé et $D=T-r=2\kappa r^2$ la part dimère, la dérivée de $T=r+2\kappa r^2$ à $T$ fixé, $0=\mathrm dr\pare{1+4\kappa r}+2r^2\,\mathrm d\kappa$, et $1+4\kappa r=(2T-r)/r$, donnent
\begin{align}
\pdp YE=\frac{n\ln10}s\,\frac{e_1\,rD}{T(2T-r)}
\et
\pdp YT=\frac{e_1\,rD}{T^2(2T-r)},
\end{align}
de signe opposé pour un dimère réduit.
Le bilan du mélange multiplié par $V_0+V$,
\begin{align}
g=\sum_i\Sigma_i\,Y_i\pare{E,T_i^\text{mél}}-\pare{V_0I_i^\text{b}+VI_i^\text{u}}
\qquad\text{avec}\qquad\Sigma_i=V_0T^\text{b}_i+VT^\text{u}_i,
\end{align}
a pour dérivées
\begin{align}
\pdp gE&=\sum_i\Sigma_i\pdp{Y_i}E,\\
\pdp gV&=\sum_i\pare{T^\text{u}_i\pare{Y_i-e_{q_i}^\text{u}}+\Sigma_i\pdp{Y_i}T\,\frac{T^\text{u}_i-T^\text{mél}_i}{V_0+V}},
\end{align}
toutes analytiques, sans différence finie (\code{electrons\_et\_derivees}).

\subsection{Les équivalences}
La burette apporte une capacité d'électrons ; le bécher offre un site par étape, de taille $T\,n_j$ et de potentiel $E^{\circ\prime}_j$, rangés par potentiel croissant pour un oxydant titrant, décroissant pour un réducteur.

\begin{center}
\titrageredoxdefinir{permanganate}{volume = 10, becher = { espece = { E0 = 0.77, C = 0.1 } }, burette = { espece = { E0 = 1.51, n = 5, H = 8, C = 0.02, forme = ox } }, volume max = 20}
\begin{tikzpicture}
\begin{axis}[width=.62\textwidth, height=5.2cm, xlabel={$V$ (mL)}, ylabel={$E$ (V/ENH)}, grid=both]
  \addplottitrage[thick]{permanganate}
  \addplottitragesauts[red, mark=*]{permanganate}
\end{axis}
\end{tikzpicture}
\captionout{\qty{10}{mL} de \ce{Fe^2+} à \qty{0.1}{mol/L} par \ce{MnO4-} à \qty{0.02}{mol/L}, $\pH=0$. Point d'inflexion à \qty{\titragesaut{permanganate}{1}}{mL}, $E=\qty{\titragesaut[E]{permanganate}{1}}{V}$ : il n'est pas au milieu du saut, $n_1\neq n_2$.}
\end{center}

% =========================================================================================
\section{La complexométrie à pH imposé}

\subsection{Le ligand libre joue le rôle du proton}
Un métal M et ses complexes successifs $\ce{ML}_j$, $j=0\dots n$, de constantes globales $\beta_j$ ($\beta_0=1$), se comportent comme un polyacide dont les \og protons\fg seraient des ligands.
Si le ligand a des acidités, seule sa forme la plus basique L complexe ; à \pH fixé, sa fraction $\alpha$ est une constante, et l'on travaille sur le ligand libre toutes formes protonées comprises, $\ell'=\ell/\alpha$, avec les constantes conditionnelles $\beta'_j=\beta_j\alpha^j$.
À $\ell'$ imposé, les fractions sont explicites,
\begin{align}
x_j=\frac{\beta'_j\ell'^j}{\sum_m\beta'_m\ell'^m},
\end{align}
calculées en échelle logarithmique comme \eqref{eq:lse}, avec $t_j=\log\beta'_j-j\,\pL$ et $\pL=-\log\ell'$ (\code{equilibre.fractions\_complexes}).
Un complexe que les tables ne donnent pas a $\log\beta=-\infty$ : sa fraction est nulle, et le calcul n'en est pas gêné.

\subsection{Le bilan du ligand}
Le ligand se répartit entre sa forme libre et les complexes : pour une solution de métaux $C_i$ et de ligand $C_L$,
\begin{align}
\rsltm{R(\pL)=\sum_iC_i\,Y_{1,i}(\pL)+\ell'-C_L}
\qquad\text{avec}\qquad Y_{1,i}=\sum_jj\,x_{j,i},
\end{align}
nul à l'équilibre.
Il est linéaire dans les concentrations de la solution, à \pL fixé : il se sépare entre le bécher et la burette comme l'électroneutralité, et $V(\pL)=-V_0\Rb(\pL)/\Ru(\pL)$, plusieurs métaux en compétition compris.
Sa dérivée,
\begin{align}
\ddp R\pL=-\ln10\pare{\sum_iC_i\,\mathrm{Var}_i+\ell'}<0,
\end{align}
avec $\mathrm{Var}_i$ écrite comme \eqref{eq:variance}, assure l'unicité de $\pL(V)$, trouvé par dichotomie sur $u=-\pL\in[-80;\log C_\text{max}+1]$ : à $\ell'=10^{-80}$, même le complexe le plus stable des tables ($\log\beta\approx40$) n'en fixe plus rien (\code{TitrageComplexometrique.pL}).

\subsection{pM et sa pente}
Le métal libre est $[\ce{M}]=T_\mathrm M(V)\,x_0$, soit
\begin{align}
\pM=-\log T_\mathrm M(V)+\log\sum_j\beta'_j\ell'^j,
\qquad
T_\mathrm M(V)=\frac{V_0C_\mathrm M^\text{b}+VC_\mathrm M^\text{u}}{V_0+V}.
\end{align}
Avec $\ddp{T_\mathrm M}V=\pare{C^\text{u}_\mathrm M-T_\mathrm M}/(V_0+V)$ et $\ddp{}\pL\log\sum_j\beta'_j\ell'^j=-Y_1$,
\begin{align}
\rsltm{\ddp\pM\pL=-\frac{C^\text{u}_\mathrm M-T_\mathrm M}{(V_0+V)\,T_\mathrm M\ln10}\ddp V\pL-Y_1}
\qquad\text{et}\qquad
\ddp\pM V=\ddp\pM\pL\Big/\ddp V\pL,
\end{align}
$\ddp V\pL$ venant de \eqref{eq:dV} appliquée au bilan du ligand.

\subsection{La courbe et les sauts}
La courbe part d'un milliardième du volume maximal : à volume nul, le bécher n'a pas de ligand libre, et \pL est infini.
Les points d'inflexion sont les minima de $\left|\ddp V\pM\right|=\left|\ddp V\pL\big/\ddp\pM\pL\right|$, cherchés en \pL et affinés comme en acide-base, à partir d'un millionième du volume maximal : avant, la courbe est plate à la précision de la machine, et ses minima n'y seraient que du bruit.

\subsection{Les équivalences}
Comme pour un polyacide, une équivalence par étape de complexation, qu'elle donne un saut ou non.
Mais les constantes successives d'un métal ne décroissent pas toujours, et un complexe peut ne jamais prédominer.
\begin{itemize}
\item Quand le ligand libre augmente, le complexe prédominant est celui qui maximise $\log\beta_j+j\log\ell$ : c'est l'enveloppe supérieure d'une famille de droites en $\log\ell$.
Par dualité, les complexes qui prédominent tour à tour sont les \gras{sommets de l'enveloppe concave} des points $(j,\log\beta_j)$, et chaque arête de cette enveloppe est une étape, qui fixe $\Delta j$ ligands avec la force $\Delta\log\beta/\Delta j$ (\code{\_sommets}, \code{\_etapes}).
\item Un ligand qui titre des métaux : les étapes de tous les métaux, rangées de la plus forte à la plus faible ; la $k$-ième équivalence apporte le ligand des $k$ premières, et deux étapes de même force n'en font qu'une, comme deux sites de même \pKa.
\item Un métal qui titre le ligand : celui-ci, d'abord en excès, est fixé dans le complexe prédominant le plus riche, puis dans les suivants : une équivalence par sommet, $V=V_0C_L/\pare{n_\text{sommet}C_\mathrm M}$.
\end{itemize}
\rmqphys{\ce{Zn}–\ce{NH3} : $\log K=\num{2,21}$, \num{2,29}, \num{2,36}, \num{2,03}, croissantes jusqu'à la troisième ; \ce{Zn(NH3)^2+} et \ce{Zn(NH3)2^2+} ne prédominent jamais, et la première étape fixe trois ligands.
Le dosage de Liebig : la source ne donne pas $\beta_1$ pour \ce{Ag}–\ce{CN-} ; \ce{Ag(CN)3^2-} puis \ce{Ag(CN)2-}, et seule la seconde donne un saut, comme la troisième acidité de l'acide phosphorique.}

\begin{center}
\titragecomplexationdefinir{calcium}{volume = 20, pH = 10, ligand pKa = {0, 1.5, 2, 2.66, 6.16, 10.24}, complexe = {metal = Ca2+, log beta = 10.7}, becher = {metal = {Ca2+, 0.01}}, burette = {ligand = 0.01}, volume max = 40}
\titragecomplexationdefinir{zinc}{volume = 20, pH = 10, ligand pKa = {0, 1.5, 2, 2.66, 6.16, 10.24}, complexe = {metal = Ca2+, log beta = 10.7}, complexe = {metal = Zn2+, log beta = 16.5}, becher = {metal = {Ca2+, 0.01}, metal = {Zn2+, 0.005}}, burette = {ligand = 0.01}, suivi = Zn2+, volume max = 40}
\begin{tikzpicture}
\begin{axis}[width=.62\textwidth, height=5.2cm, xlabel={$V$ (mL)}, ylabel=\pM, grid=both]
  \addplottitrage[thick]{calcium}
  \addplottitragesauts[black, mark=*]{calcium}
  \addplottitrage[thick]{zinc}
  \addplottitragesauts[black, mark=*]{zinc}
\end{axis}
\end{tikzpicture}
\captionout{\qty{20}{mL} de \ce{Ca^2+} à \qty{0.01}{mol/L} par l'EDTA à \qty{0.01}{mol/L}, à $\pH=10$ : pCa \num{\titragepM{calcium}{20}} à l'équivalence. Avec \ce{Zn^2+} à \qty{0.005}{mol/L}, pZn : équivalences \num{\titrageequivalence{zinc}{1}} et \qty{\titrageequivalence{zinc}{2}}{mL}, le zinc d'abord ($\log\beta=\num{16,5}$).}
\end{center}

% =========================================================================================
\section{La précipitation à pH imposé}

\subsection{Les solides présents}
Un ion commun P (\ce{Ag+}) forme avec des partenaires $\mathrm X_i$ des solides $\mathrm P_\nu\mathrm X_\mu$, $K_s=a^\nu x^\mu$ avec $a=[\mathrm P]$.
À $a$ imposé, un solide présent fixe son partenaire, toutes formes protonées comprises,
\begin{align}
x_\text{sat}=\pare{K_s/a^\nu}^{1/\mu}\big/\alpha,
\end{align}
et un solide absent laisse le partenaire en solution.
En notant $r=\nu/\mu$, pour un ensemble $S$ de solides présents,
\begin{align}
\rsltm{R_S(\pP)=a+\sum_{i\in S}r_i\pare{C_{\mathrm X_i}-x_{\text{sat},i}}-C_\mathrm P}
\end{align}
est encore linéaire dans les concentrations de la solution : chaque ensemble a son volume explicite, $V_S(\pP)=-V_0R_{S,\text{bécher}}/R_{S,\text{burette}}$.

\subsection{L'ensemble cohérent}
On essaie les $2^n$ ensembles, et l'on retient celui qui est \gras{cohérent} au volume trouvé : chaque solide présent a une quantité positive, $T_{\mathrm X_i}(V)\ge x_{\text{sat},i}$, chaque absent n'est pas sursaturé, $T_{\mathrm X_i}(V)\le x_{\text{sat},i}$, à $10^{-12}$ près en relatif ; un volume de $-10^{-12}V_0$, au départ, est un zéro arrondi.
À volume donné, l'ensemble ne se choisit pas : la fonction
\begin{align}
F_V(a)=a+\sum_ir_i\max\pare{0,T_{\mathrm X_i}(V)-x_{\text{sat},i}(a)}-T_\mathrm P(V)
\end{align}
est strictement croissante en $a$ — $x_\text{sat}$ décroît avec $a$ —, son zéro est unique, et les solides présents en découlent ; c'est la dichotomie de \code{TitrageParPrecipitation.p}, sur $\log a\in[-60;\log C_\text{max}+1]$.

\subsection{La pente}
La dérivée de $R_S$ ne dépend que de $a$ et des solides présents,
\begin{align}
\ddp{R_S}\pP=-\ln10\pare{a+\sum_{i\in S}r_i^2x_{\text{sat},i}},
\end{align}
la même pour les deux solutions ; \eqref{eq:dV} donne alors $\ddp V\pP=-(V_0+V)\,R_S'/R_{S,\text{burette}}$.

\subsection{Les événements}
Là où un solide apparaît ou disparaît, la pente saute : la courbe a un \gras{point anguleux}, un événement, distinct des points d'inflexion.
On les repère où l'ensemble change sur une grille de 4000 \pP, et on les affine par dichotomie sur le signe de la sursaturation $T_\mathrm X(V)-x_\text{sat}(a)$.
Les points d'inflexion ne se cherchent qu'à l'intérieur des morceaux lisses : les mêmes solides de part et d'autre.

\subsection{Les équivalences}
L'ion commun titre les partenaires dans l'ordre où ils précipitent dans le bécher, par le \pP où chacun commencerait à précipiter, $(\mu(\log C+\log\alpha)+\pKs)/\nu$ ; un partenaire titre l'ion commun qui reste une fois les partenaires du bécher précipités (Charpentier-Volhard).

\begin{center}
\titrageprecipitationdefinir{mohr}{volume = 20, pH = 7.5, precipite = {partenaire = Cl-, pKs = 9.75}, precipite = {partenaire = CrO4^2-, pKs = 11.95, nu = 2, pKa = 6.5}, becher = {partenaire = {Cl-, 0.05}, partenaire = {CrO4^2-, 0.005}}, burette = {commun = 0.05}, volume max = 30}
\titrageprecipitationdefinir{volhard}{volume = 20, precipite = {partenaire = Cl-, pKs = 9.75}, precipite = {partenaire = SCN-, pKs = 12}, becher = {commun = 0.1, partenaire = {Cl-, 0.05}}, burette = {partenaire = {SCN-, 0.1}}, volume max = 25}
\begin{tikzpicture}
\begin{axis}[width=.62\textwidth, height=5.2cm, xlabel={$V$ (mL)}, ylabel={pAg}, grid=both]
  \addplottitrage[thick]{mohr}
  \addplottitrageevenements[black, mark=diamond*]{mohr}
  \addplottitrage[thick]{volhard}
  \addplottitrageevenements[black, mark=diamond*]{volhard}
\end{axis}
\end{tikzpicture}
\captionout{Mohr, $\pH=\num{7,5}$ : le chromate d'argent apparaît à \qty{\titrageevenement{mohr}{2}}{mL}. Charpentier-Volhard : le thiocyanate d'argent apparaît dès la première goutte, le chlorure d'argent se redissout à \qty{\titrageevenement{volhard}{2}}{mL}. Losanges : les événements.}
\end{center}

% =========================================================================================
\section{Le pH calculé}

\subsection{Les composantes}
Sans tampon, le pH devient une inconnue.
On reprend la formulation des logiciels de spéciation (la méthode du tableau de Morel) : les inconnues sont les concentrations libres de quelques \gras{composantes} — \ce{H+}, le ligand sous sa forme la plus basique, l'ion commun, les métaux, les partenaires —, en logarithme, $u_c=\log c$.
Toute espèce s'en déduit, $c_s=K_s\,10^{\nu_s\cdot u}$ : $\ce{H_kL}$ a $\nu_\mathrm H=k$, $\nu_\mathrm L=1$ ; \ce{HO-} a $\nu_\mathrm H=-1$ et $K=\Ke$.
Chaque composante a son bilan, $\sum_s\nu_{sc}c_s=T_c$, et chaque total est linéaire dans les deux solutions.
Le bilan de H est l'électroneutralité écrite en protons : un acide fort compte dans $T_\mathrm H$, une base forte en négatif, \ce{Na2H2Y} apporte $2C$ à $T_\mathrm H$.

\subsection{L'unicité : la convexité}
Les bilans sont le gradient de
\begin{align}
\Phi(u)=\frac1{\ln10}\sum_sc_s(u)-\sum_cT_cu_c,
\qquad
\pdp\Phi{u_c}=\sum_s\nu_{sc}c_s-T_c,
\end{align}
dont la hessienne, $\ln10\sum_sc_s\,\nu_s\nu_s^\top$, est définie positive dès que les composantes sont indépendantes : $\Phi$ est strictement convexe, et l'équilibre à volume donné est \gras{unique}.
Un solide ajoute une contrainte linéaire, $\nu_p\cdot u\le\log K_{s,p}$ : le domaine reste convexe, la solution unique, et la quantité de solide est le multiplicateur de la contrainte.

\rmqphys{La conséquence qui fonde l'algorithme : minimiser une fonction convexe sur une partie des variables laisse une fonction convexe des autres, dont la dérivée est le bilan restant.
\gras{Le bilan des protons, évalué après avoir résolu les autres bilans à \pH fixé, est donc monotone en \pH.}
D'où deux dichotomies imbriquées, chacune sur une fonction monotone, sans Newton ni point de départ (\code{\_dichotomie}).}

Une composante est \gras{éliminable} quand son bilan se résout explicitement une fois les autres fixées : un métal, à $\ell$ connu, par ses fractions ; un partenaire, à $a$ et \pH connus, par $x'=\min\pare{T_\mathrm X,x_\text{sat}}$ ; un métal qui précipite en hydroxyde, par la borne qu'impose \pH.
Restent au plus deux composantes principales : le \pH, à l'extérieur, et $u$ — le ligand ou l'ion commun —, à l'intérieur.

\subsection{La complexométrie}
Avec $\ell$ le ligand basique libre, $P=\sum_kK_kh^k$ (les $K_k$ cumulés de ses formes protonées), $P_1=\sum_kk\,K_kh^k$ et $P_2=\sum_kk^2K_kh^k$, les deux bilans du mélange sont
\begin{align}
g_L&=\sum_mT_mY_{1,m}(u)+\ell P-T_L,\\
g_H&=R_\text{ab}(\pH)+k_0T_L-\ell P_1,
\end{align}
où $R_\text{ab}$ est le résidu \eqref{eq:R} de la partie acido-basique des solutions (acides et bases forts, un tampon décrit comme un couple, l'eau), $k_0$ les protons du ligand tel qu'on l'introduit, et $\ell P_1$ ceux que portent ses formes libres ; le ligand fixé dans un complexe n'en porte pas.
La jacobienne, en $(\pH,u)$ :
\begin{align}
\pdp{g_H}\pH=R_\text{ab}'+\ln10\,\ell P_2,\quad
\pdp{g_H}u=\pdp{g_L}\pH=-\ln10\,\ell P_1,\quad
\pdp{g_L}u=\ln10\pare{\sum_mT_m\mathrm{Var}_m+\ell P},
\end{align}
symétrique, comme le veut la convexité (\code{complexation\_libre.py}).

\subsection{La précipitation et les hydroxydes}
Avec $a$ l'ion commun libre, $x'_i=x_{\text{sat},i}$ pour un solide présent et $T_{\mathrm X_i}$ sinon, $m_i$ le nombre moyen de protons du partenaire en solution :
\begin{align}
g_P&=a+\sum_{i\in S}r_i\pare{T_{\mathrm X_i}-x'_i}-T_P,\\
g_H&=R_\text{ab}(\pH)-\sum_ix'_im_i,
\end{align}
un partenaire s'introduisant sous sa forme la plus basique.
Un \gras{hydroxyde} $\mathrm P_\nu(\mathrm{OH})_\mu$ — un précipité dont le partenaire est \ce{HO-} — borne l'ion commun : $a\le a_\text{sat}(h)$, $K_s=a^\nu(\Ke/h)^\mu$.
Présent, il fixe $u=\pare{-\pKs-\mu(\pH-\pKe)}/\nu$, et sa quantité $s=-g_P/\nu$ libère $\mu s$ protons : l'équation du \pH devient $g_H-(\mu/\nu)\,g_P=0$, à une seule inconnue (\code{precipitation\_libre.py}).

\subsection{Newton, et son critère d'arrêt}
Pour suivre la courbe, on n'emploie la dichotomie imbriquée qu'en reprise : chaque point part de l'état d'un voisin (\pH, $u$, solides), et Newton, à jacobienne exacte et à ensemble de solides fixé, y converge en quelques pas.
\begin{itemize}
\item Les pas sont bornés à deux unités, pour qu'un départ lointain ne sorte pas du domaine.
\item Près d'une équivalence, le bilan du ligand est mal conditionné : $\pdp{g_L}u$ y est petit, et le bruit d'arrondi de $g_L$ ($\approx10^{-18}$ sur des totaux de $10^{-2}$) laisse $u$ indéterminé à quelques $10^{-13}$ près, pour Newton comme pour la dichotomie.
Le critère d'arrêt en tient compte : le pas tombe sous $10^{-13}$, ou cesse de décroître alors qu'il est déjà sous $10^{-8}$.
\item Un point où Newton n'a pas convergé, ou dont l'ensemble de solides ne convient plus, est repris par la dichotomie, qui est sûre.
\end{itemize}

\subsection{Les dérivées}
À volume fixé, le théorème des fonctions implicites donne la pente exactement :
\begin{align}
J\begin{pmatrix}\ddp\pH V\\[2pt]\ddp uV\end{pmatrix}=-\pdp gV=-\frac1{V_0+V}\begin{pmatrix}g_H^\text{burette}\\ g_{L}^\text{burette}\end{pmatrix},
\end{align}
où $g^\text{burette}$ est le bilan écrit avec les seules concentrations de la burette, termes libres compris : sur la courbe, $g=0$, et $\pdp{}V\frac{V_0g^\text{b}+Vg^\text{u}}{V_0+V}=\frac{g^\text{u}-g}{V_0+V}$.

\subsection{La courbe paramétrée par le volume}
Le volume n'est plus une fonction explicite du paramètre, et le \pH n'est pas toujours monotone le long du titrage : la courbe est paramétrée par $V$ lui-même.
On part de 200 volumes réguliers, puis on ajoute le milieu de chaque segment plus long que $1/n$ dans le plan normalisé, jusqu'à ce qu'il n'en reste plus ; on rééchantillonne enfin à longueur d'arc constante (\code{libre.py}).
Les points d'inflexion sont les minima de $\left|\ddp V\pM\right|$ en $V$, à l'intérieur des morceaux lisses, affinés par \code{minimum\_lisse} avec un pas proportionnel à la largeur locale de la grille raffinée : avec des dérivées exactes à $10^{-15}$ près, la position est juste à environ $10^{-12}$ fois cette largeur.
Les événements sont affinés par dichotomie en volume, sur le prolongement du morceau qui les précède.

\begin{center}
\titragecomplexationdefinir{sanstampon}{volume = 20, ligand pKa = {0, 1.5, 2, 2.66, 6.16, 10.24}, complexe = {metal = Ca2+, log beta = 10.7}, becher = {metal = {Ca2+, 0.01}}, burette = {ligand = 0.01, protons ligand = 2}, volume max = 40}
\titragecomplexationdefinir{ammoniacal}{volume = 20, ligand pKa = {0, 1.5, 2, 2.66, 6.16, 10.24}, complexe = {metal = Ca2+, log beta = 10.7}, becher = {metal = {Ca2+, 0.01}, espece = {pKa = 9.25, C = 0.1}, espece = {pKa = 9.25, C = 0.1, forme = 1}}, burette = {ligand = 0.01, protons ligand = 2, espece = {pKa = 9.25, C = 0.1}, espece = {pKa = 9.25, C = 0.1, forme = 1}}, volume max = 40}
\begin{tikzpicture}
\begin{axis}[width=.62\textwidth, height=5.2cm, xlabel={$V$ (mL)}, ylabel={pCa}, grid=both]
  \addplottitrage[thick]{sanstampon}
  \addplottitrage[thick]{ammoniacal}
  \addplottitragesauts[black, mark=*]{ammoniacal}
\end{axis}
\end{tikzpicture}
\captionout{\ce{Ca^2+} à \qty{0.01}{mol/L} par \ce{Na2H2Y} à \qty{0.01}{mol/L}, pH calculé. Sans tampon, l'EDTA libère ses protons : \pH \num{\titragepH{sanstampon}{10}} à \qty{10}{mL}, plus de saut. Avec le tampon \ce{NH4+}/\ce{NH3} à \qty{0.1}{mol/L} décrit comme une solution, le \pH dérive de \num{\titragepH{ammoniacal}{0}} à \num{\titragepH{ammoniacal}{20}}, et le saut reste, à \qty{\titragesaut{ammoniacal}{1}}{mL}.}
\end{center}

% =========================================================================================
\section{Les trois implémentations}

\begin{center}
\begin{tabularx}{\linewidth}{@{}>{\raggedright\arraybackslash}X>{\raggedright\arraybackslash}p{.27\linewidth}>{\raggedright\arraybackslash}p{.22\linewidth}@{}}
\toprule
& Python et Lua & pgf (pdf\LaTeX)\\
\midrule
courbes acide-base, conductimétrie, rédox monomère & oui & oui \\
couples dimériques, sauts, équivalences, tangentes & oui & non \\
complexométrie, précipitation, pH calculé & oui & non \\
précision des courbes & double précision & \code{fpu}, 4 à 5 chiffres \\
précision des nombres & double précision & \code{l3fp}, 16 chiffres \\
\bottomrule
\end{tabularx}
\end{center}
\begin{itemize}
\item Le moteur Python est vectorisé : les dichotomies avancent sur des tableaux de points d'un coup.
\item Le moteur Lua suit le même algorithme, point par point, en Lua 5.3 pur ; il tourne sous Lua\TeX{} comme sous \code{texlua}, ce qui permet de le comparer au moteur Python hors de \TeX.
Pour le pH calculé, où une dichotomie imbriquée coûte quelques dizaines de millisecondes par point, la courbe se construit par continuation : Newton depuis le point précédent, la dichotomie en reprise.
\item Le moteur pgf construit, à la définition du titrage, l'expression de $V(\pH)$ et la confie à \code{pgfplots} en courbe paramétrée ; la même expression, évaluée par \code{l3fp}, donne les nombres explicites, et ceux qu'une dichotomie de 54 coupes trouve ($\num{2e-15}$ en \pH).
\end{itemize}

% =========================================================================================
\section{Les vérifications}

\subsection{Le principe}
Chaque calcul du paquet est confronté à un calcul qui ne lui emprunte rien : ni son code, ni, autant que possible, sa formulation.
\begin{itemize}
\item \gras{La dilution} : la force brute dilue à la main chaque concentration, écrit chaque forme par la formule des manuels, avec ses contre-ions, et résout l'électroneutralité du mélange par 200 coupes de dichotomie en double précision ; elle ne connaît pas la séparation \eqref{eq:V}.
\item \gras{Les points d'inflexion} : mpmath, à 60 chiffres, écrit la courbe directe en multiprécision et cherche le zéro de sa dérivée seconde (\code{mp.findroot}), depuis le minimum de $|V'|$ sur une grille (\code{tests/reference\_mpmath.py}).
\item \gras{La complexométrie et la précipitation} : mpmath, à 50 chiffres, dilue à la main et résout les bilans du mélange par dichotomie sur le ligand libre ou sur l'ion commun, chaque solide présent ou non selon qu'il est sursaturé ; il ne connaît ni le paramétrage par \pL ou \pP, ni l'énumération des ensembles.
\item \gras{Le pH calculé} : la référence écrit l'électroneutralité, avec les contre-ions écrits à la main (\ce{Cl-} de \ce{CaCl2}, \ce{Na+} de \ce{Na2H2Y}, \ce{K+} de \ce{KCN}…), là où le paquet écrit le bilan des protons ; elle trouve un départ par dichotomie en double précision, puis résout le système de l'ensemble de solides trouvé par Newton à 40 chiffres (\code{mp.findroot}).
\item \gras{Les moteurs Lua et pgf} : leurs nombres, comparés à ceux du paquet Python sur les mêmes titrages.
\end{itemize}
Les tables suivantes donnent, pour chaque cas, le plus grand écart mesuré, et non la tolérance des tests.
Elles sont écrites par \code{docs/technique/verifications.py}, qui reprend les références des tests (\code{tests/}).

\subsection{La dilution}
\begin{center}
\captionof{table}{Le \pH (ou le potentiel, en volts) du paquet face à la force brute qui dilue à la main, au même volume versé.}\label{tab:dilution}
% écrit par docs/technique/verifications.py : ne pas éditer
\begin{tabular}{lrr}
\toprule
Titrage & volumes & écart maximal (pH ou V) \\
\midrule
acide fort par base forte & 12 & \num{2.3e-11} \\
acide faible par base forte & 12 & \num{4.1e-13} \\
acide faible dilué par base concentrée & 12 & 0 \\
acide faible concentré par base diluée & 12 & \num{8.9e-16} \\
triacide par base forte & 12 & \num{3.6e-14} \\
base faible par acide fort & 12 & \num{9.4e-14} \\
sel de diacide par acide fort & 12 & \num{3.4e-14} \\
mélange par base forte & 12 & \num{7.1e-15} \\
acide faible par base faible & 12 & \num{1.1e-14} \\
tampon par tampon & 12 & 0 \\
\ce{Fe2+} par \ce{Ce4+} ($E$) & 6 & 0 \\
\ce{Fe2+} par \ce{Cr2O7^2-}, dimère ($E$) & 6 & 0 \\
\bottomrule
\end{tabular}

\end{center}
\rmqcons{L'écart de \num{2.3e-11} de l'acide fort est à l'équivalence, où $\ddp\pH V$ est très grande : un écart relatif de $10^{-15}$ sur le volume y devient $10^{-11}$ en \pH. Ce n'est pas une erreur du calcul, c'est le conditionnement du problème.}

\subsection{Les points d'inflexion}
\begin{center}
\captionof{table}{Le point d'inflexion, en \pH ou en volts/ENH : mpmath à 60 chiffres face au paquet. Les volumes valent \qty{10}{mL} pour \qty{10}{mL} de bécher.}\label{tab:inflexions}
% écrit par docs/technique/verifications.py : ne pas éditer
\begin{tabular}{lrrr}
\toprule
Titrage & mpmath & paquet & écart \\
\midrule
\ce{CH3COOH} par \ce{NaOH} & \num{8.729420699201} & \num{8.7294206993005} & \num{1.0e-10} \\
\ce{HCl} par \ce{NaOH} & \num{6.999998262822} & \num{6.9999982628470} & \num{2.5e-11} \\
\ce{NH3} par \ce{HCl} & \num{5.275575942902} & \num{5.2755759429165} & \num{1.4e-11} \\
\ce{Fe2+} par \ce{Ce4+} (V/ENH) & \num{1.244999999519} & \num{1.2449999995252} & \num{5.8e-12} \\
\ce{Fe2+} par \ce{MnO4-} (V/ENH) & \num{1.37288297818} & \num{1.3728829781858} & \num{5.6e-12} \\
\ce{Fe2+} par \ce{Cr2O7^2-} (V/ENH) & \num{1.249330805271} & \num{1.2493308052717} & \num{1.0e-12} \\
\bottomrule
\end{tabular}

\end{center}

\subsection{La complexométrie et la précipitation à pH imposé}
\begin{center}
\captionof{table}{pM du paquet face à mpmath à 50 chiffres (\code{tests/test\_complexation.py}).}\label{tab:complexation}
% écrit par docs/technique/verifications.py : ne pas éditer
\begin{tabular}{lrr}
\toprule
Titrage & volumes & écart maximal sur pM \\
\midrule
\ce{Ca}-EDTA, pH 10 & 10 & \num{4.8e-13} \\
\ce{Ca}-EDTA, pH 6 & 10 & \num{2.7e-15} \\
\ce{Ca} et \ce{Zn} en compétition, p\ce{Zn} & 10 & \num{2.0e-14} \\
\ce{Ca} et \ce{Zn} en compétition, p\ce{Ca} & 10 & \num{8.9e-16} \\
complexes successifs \ce{Cu}-\ce{NH3}, ligand titrant & 10 & \num{1.8e-15} \\
complexes successifs \ce{Cu}-\ce{NH3}, métal titrant & 9 & \num{2.7e-15} \\
ligand sans acidité & 10 & \num{9.8e-15} \\
\bottomrule
\end{tabular}


\bigskip
\captionof{table}{pP du paquet face à mpmath à 50 chiffres (\code{tests/test\_precipitation.py}).}\label{tab:precipitation}
% écrit par docs/technique/verifications.py : ne pas éditer
\begin{tabular}{lrr}
\toprule
Titrage & volumes & écart maximal sur pP \\
\midrule
chlorure par l'argent & 11 & \num{1.1e-13} \\
chlorure et iodure & 11 & \num{1.7e-13} \\
Mohr, pH 7,5 & 11 & \num{7.3e-14} \\
Charpentier-Volhard & 11 & \num{9.9e-14} \\
argent par le chlorure & 11 & \num{1.1e-13} \\
\bottomrule
\end{tabular}

\end{center}

\subsection{Le pH calculé}
\begin{center}
\captionof{table}{Complexométrie à pH calculé : pH et pM du paquet face à la référence par l'électroneutralité, Newton à 40 chiffres (\code{tests/test\_complexation\_libre.py}).}\label{tab:complexation-libre}
% écrit par docs/technique/verifications.py : ne pas éditer
\begin{tabular}{lrrr}
\toprule
Titrage & volumes & écart sur pH & écart sur pM \\
\midrule
\ce{Ca} par \ce{Na2H2Y}, sans tampon & 9 & \num{4.9e-15} & \num{1.3e-15} \\
\ce{Ca} par \ce{Na2H2Y}, tampon ammoniacal en solution & 9 & \num{1.8e-15} & \num{2.3e-14} \\
\ce{Zn} par \ce{Na2H2Y}, avec de l'acide fort & 9 & \num{4.0e-15} & \num{9.8e-15} \\
\ce{Cu} par l'ammoniac, sans tampon & 9 & \num{1.8e-15} & \num{1.8e-15} \\
Liebig, avec la protonation de \ce{HCN} & 9 & \num{1.1e-12} & \num{4.5e-12} \\
\bottomrule
\end{tabular}


\bigskip
\captionof{table}{Précipitation à pH calculé, hydroxydes compris (\code{tests/test\_precipitation\_libre.py}).}\label{tab:precipitation-libre}
% écrit par docs/technique/verifications.py : ne pas éditer
\begin{tabular}{lrrr}
\toprule
Titrage & volumes & écart sur pH & écart sur pP \\
\midrule
Mohr sans tampon & 8 & \num{2.5e-12} & \num{4.4e-14} \\
Mohr en milieu acide nitrique & 8 & \num{3.6e-15} & \num{4.4e-14} \\
Charpentier-Volhard en milieu nitrique & 7 & \num{2.2e-16} & \num{1.3e-13} \\
\ce{Mg(OH)2} par la soude & 7 & \num{1.8e-15} & \num{4.4e-15} \\
\ce{Mg(OH)2} après un acide fort & 9 & \num{2.7e-12} & \num{2.7e-15} \\
\ce{Ag2O} par la soude & 6 & \num{1.8e-15} & \num{1.8e-15} \\
\bottomrule
\end{tabular}

\end{center}
\rmqcons{Les deux limites du niveau imposé sont vérifiées à part : un ligand ou un partenaire sans acidité, dont le pH se calcule sans compter, donne le même pM ou pP qu'à pH imposé à $10^{-12}$ près ; un tampon de plus en plus concentré rapproche le résultat de celui du pH imposé, l'écart diminuant d'un facteur 5 au moins à chaque décade de concentration.}

\subsection{Le moteur Lua}
\begin{center}
\captionof{table}{Le moteur Lua face au moteur Python, sur les mêmes titrages (\code{tests/test\_lua.py}) : la courbe (pH, $E$, pM ou pP, à des volumes donnés), puis la position relative et la grandeur des points d'inflexion.}\label{tab:lua}
\resizebox{\linewidth}{!}{% écrit par docs/technique/verifications.py : ne pas éditer
\begin{tabular}{llrrr}
\toprule
Famille & Titrage & écart sur la courbe & sauts : $V$ (relatif) & sauts : grandeur \\
\midrule
acide-base & \ce{CH3COOH} par \ce{NaOH} & 0 & 0 & \num{2.1e-14} \\
acide-base & \ce{H3PO4} par \ce{NaOH} & 0 & 0 & 0 \\
acide-base & \ce{NH3} par \ce{HCl} & 0 & \num{7.3e-15} & \num{1.6e-11} \\
acide-base & \ce{HCl} et \ce{CH3COOH} par \ce{NH3} & 0 & \num{1.5e-12} & \num{1.5e-11} \\
rédox & \ce{Fe2+} par \ce{Ce4+} & 0 & \num{1.8e-16} & \num{2.2e-15} \\
rédox & \ce{Fe2+} par \ce{Cr2O7^2-} & 0 & 0 & \num{2.3e-13} \\
rédox & \ce{I2} par \ce{S2O3^2-} & 0 & \num{3.0e-15} & \num{6.0e-12} \\
complexation & \ce{Ca}-EDTA & 0 & 0 & 0 \\
complexation & \ce{Ca} et \ce{Zn} & 0 & 0 & 0 \\
complexation & \ce{Cu}-\ce{NH3}, métal titrant & 0 & 0 & 0 \\
précipitation & Mohr & 0 & 0 & 0 \\
précipitation & iodure et chlorure & 0 & 0 & \num{7.1e-15} \\
précipitation & Charpentier-Volhard & 0 & 0 & 0 \\
complexation, pH calculé & \ce{Ca} par \ce{Na2H2Y} & 0 & 0 & 0 \\
complexation, pH calculé & \ce{Ca}, tampon en solution & 0 & \num{9.8e-14} & \num{9.3e-11} \\
complexation, pH calculé & \ce{Cu} par \ce{NH3} & 0 & \num{1.4e-11} & \num{5.6e-11} \\
complexation, pH calculé & Liebig & 0 & \num{7.1e-15} & \num{4.8e-10} \\
précipitation, pH calculé & Mohr sans tampon & 0 & \num{1.8e-14} & \num{6.6e-13} \\
précipitation, pH calculé & Mohr en milieu acide & 0 & \num{8.2e-14} & \num{3.8e-12} \\
précipitation, pH calculé & Volhard en milieu nitrique & 0 & \num{7.1e-15} & \num{3.7e-12} \\
précipitation, pH calculé & \ce{Mg(OH)2} après un acide & 0 & 0 & 0 \\
précipitation, pH calculé & \ce{Ag2O} & 0 & 0 & \num{8.9e-16} \\
\bottomrule
\end{tabular}
}
\end{center}
\rmqphys{Les courbes sont identiques au bit près, ou à quelques $10^{-14}$ : les deux moteurs font les mêmes opérations dans le même ordre.
Les points d'inflexion diffèrent davantage en grandeur qu'en volume : un saut amplifie l'écart, comme pour la dilution.}

\subsection{Le moteur pgf}
Sous pdf\LaTeX, \code{tests/test\_latex.py} compile \code{latex/tests/titrages-test-nombres.tex} et compare ses nombres au paquet Python : ceux d'\code{l3fp} à $10^{-9}$ près ; les courbes, évaluées par l'arithmétique de \code{pgfplots} (\code{fpu}), à $\num{2e-3}$ près en relatif, moins que l'épaisseur du trait.

\subsection{La suite de tests}
La suite compte 740 tests ; outre ces comparaisons, elle vérifie les pentes contre des différences finies, les équivalences, les doctests de la documentation, et les deux paquets \LaTeX{} compilés.
Elle se rejoue par
\begin{center}
\code{uv run pytest} \quad et \quad \code{cd latex/tests \&\& texlua test\_titrages.lua}
\end{center}
et les tables de ce document par \code{uv run python docs/technique/verifications.py}.

% =========================================================================================
\section{Limites}

\begin{itemize}
\item Les solutions sont idéales ; les conductivités molaires sont les valeurs limites, qui surestiment $\sigma$ dès quelques \unit{mmol/L}.
\item Un seul ligand par titrage complexométrique : un tampon qui complexerait aussi le métal, ou un indicateur métallochrome, en serait un second, et demanderait trois niveaux de dichotomie.
Complexation et précipitation ne se combinent pas dans un même titrage.
\item La courbe d'une précipitation à pH imposé est paramétrée par \pP : elle suppose \pP monotone le long du titrage, ce qui vaut dans les cas traités et se vérifie par la référence.
\item Les constantes sont celles qu'on donne, ou celles de la base, tirées des tables à \qty{25}{\degreeCelsius}.
\end{itemize}

\section*{Références}
\begin{itemize}
\item H. Kalka, \emph{Polyprotic Acids and Beyond — An Algebraic Approach}, Chemistry 2021, 3, 454-508, \href{https://doi.org/10.3390/chemistry3020034}{doi:10.3390/chemistry3020034}.
\item T. Brière, \emph{Une méthode simple de résolution des problèmes de pH et de titrages acido-basiques sans approximations}, BUP 841, 2002.
\item Y. Cortial, \emph{Courbe de dosage acidobasique et méthode des tangentes}, BUP 734, 1991.
\item F. M. M. Morel, J. G. Hering, \emph{Principles and Applications of Aquatic Chemistry}, Wiley, 1993 : la méthode du tableau et les composantes.
\item D. Harvey, \emph{Analytical Chemistry 2.1}, LibreTexts, § 16.10 à 16.12, d'après A. E. Martell, R. M. Smith, \emph{Critical Stability Constants}, Plenum, 1976.
\end{itemize}

\end{document}
